\section{Технический регламент и инфраструктура}

\begin{frame}{Технический регламент туров}
  \vspace{0.2em}
  \begin{columns}[T]
    \begin{column}{0.49\textwidth}
      \begin{block}{Среда выполнения}
        \begin{itemize}\setlength\itemsep{0.25em}
          \item доступ к \textbf{АТС} --- обеспечивает Оргкомитет;
          \item среда разработки \textbf{Jupyter Notebook} ---
                на выделенной инфраструктуре либо на предоставленных
                участникам устройствах;
          \item основной язык --- \textbf{Python};
          \item перечни разрешённых библиотек, фреймворков,
                предобученных моделей и наборов данных публикуются
                на портале \textbf{не позднее чем за 30 дней}.
        \end{itemize}
      \end{block}
    \end{column}
    \begin{column}{0.49\textwidth}
      \begin{block}{Проверка и вопросы по ходу тура}
        \begin{itemize}\setlength\itemsep{0.25em}
          \item автоматическая проверка по заранее определённым
                метрикам качества с последующей верификацией Жюри;
          \item вопросы по условию, данным и проверяющим тестам ---
                через АТС в течение тура.
        \end{itemize}
      \end{block}
      \begin{block}{Изолированная генеративная модель}
        Оргкомитет вправе предоставить участникам доступ к изолированной
        генеративной модели \textbf{до 50 млрд параметров}.
      \end{block}
    \end{column}
  \end{columns}
\end{frame}
